\documentclass[10pt,twoside]{article}\usepackage[]{graphicx}\usepackage[dvipsnames,svgnames,table]{xcolor}
% maxwidth is the original width if it is less than linewidth
% otherwise use linewidth (to make sure the graphics do not exceed the margin)
\makeatletter
\def\maxwidth{ %
  \ifdim\Gin@nat@width>\linewidth
    \linewidth
  \else
    \Gin@nat@width
  \fi
}
\makeatother

\definecolor{fgcolor}{rgb}{0.345, 0.345, 0.345}
\newcommand{\hlnum}[1]{\textcolor[rgb]{0.686,0.059,0.569}{#1}}%
\newcommand{\hlsng}[1]{\textcolor[rgb]{0.192,0.494,0.8}{#1}}%
\newcommand{\hlcom}[1]{\textcolor[rgb]{0.678,0.584,0.686}{\textit{#1}}}%
\newcommand{\hlopt}[1]{\textcolor[rgb]{0,0,0}{#1}}%
\newcommand{\hldef}[1]{\textcolor[rgb]{0.345,0.345,0.345}{#1}}%
\newcommand{\hlkwa}[1]{\textcolor[rgb]{0.161,0.373,0.58}{\textbf{#1}}}%
\newcommand{\hlkwb}[1]{\textcolor[rgb]{0.69,0.353,0.396}{#1}}%
\newcommand{\hlkwc}[1]{\textcolor[rgb]{0.333,0.667,0.333}{#1}}%
\newcommand{\hlkwd}[1]{\textcolor[rgb]{0.737,0.353,0.396}{\textbf{#1}}}%
\let\hlipl\hlkwb

\usepackage{framed}
\makeatletter
\newenvironment{kframe}{%
 \def\at@end@of@kframe{}%
 \ifinner\ifhmode%
  \def\at@end@of@kframe{\end{minipage}}%
  \begin{minipage}{\columnwidth}%
 \fi\fi%
 \def\FrameCommand##1{\hskip\@totalleftmargin \hskip-\fboxsep
 \colorbox{shadecolor}{##1}\hskip-\fboxsep
     % There is no \\@totalrightmargin, so:
     \hskip-\linewidth \hskip-\@totalleftmargin \hskip\columnwidth}%
 \MakeFramed {\advance\hsize-\width
   \@totalleftmargin\z@ \linewidth\hsize
   \@setminipage}}%
 {\par\unskip\endMakeFramed%
 \at@end@of@kframe}
\makeatother

\definecolor{shadecolor}{rgb}{.97, .97, .97}
\definecolor{messagecolor}{rgb}{0, 0, 0}
\definecolor{warningcolor}{rgb}{1, 0, 1}
\definecolor{errorcolor}{rgb}{1, 0, 0}
\newenvironment{knitrout}{}{} % an empty environment to be redefined in TeX

\usepackage{alltt}
\usepackage[marginparsep=1em]{geometry}
\geometry{lmargin=1.0in,rmargin=1.0in, bmargin=1.2in,  tmargin=1.2in}
\usepackage[dvipsnames,svgnames,table]{xcolor}
\usepackage{graphicx}
\usepackage{amssymb}
\usepackage{epstopdf}
\usepackage{verbatim}
\usepackage{enumerate}
\usepackage{bm}
\usepackage{amsthm}
\usepackage{float}
\usepackage{amsmath}
\usepackage{fancyhdr}
\usepackage{hyperref}
\usepackage{mathtools}

%%%%  SHORTCUT COMMANDS  %%%%
\newcommand{\ds}{\displaystyle}
\newcommand{\Z}{\mathbb{Z}}
\newcommand{\T}{\mathcal{T}}
\newcommand{\arc}{\rightarrow}
\newcommand{\R}{\mathbb{R}}
\newcommand{\RP}{\mathbb{R}(+)}
\newcommand{\Rs}{\mathbb{R}^{**}}
\newcommand{\C}{\mathbb{C}}
\newcommand{\E}{\mathbb{E}}
\newcommand{\B}{\mathcal{B}}
\newcommand{\PX}{\mathcal{P}(X)}
\newcommand*\rot{\rotatebox{90}}
\newcommand*\OK{\ding{51}}
\newcommand{\N}{\mathbb{N}}
\newcommand{\Q}{\mathbb{Q}}
\newcommand{\Answer}{\vspace{2mm}\textbf{\underline{Answer}}\\}
\newcolumntype{L}{>{$}l<{$}}
\newcolumntype{C}{>{$}c<{$}}
\newcommand{\GL}{\text{GL}}
\newcommand{\SL}{\text{SL}}

\newcommand{\stirling}[2]{\genfrac{\{}{\}}{0pt}{}{#1}{#2}}

\pagestyle{fancy}
\fancyhf{}
\renewcommand{\sectionmark}[1]{\markright{\thesection.\ #1}}
\lhead{\fancyplain{}{}}
\fancyhead[RE,RO]{Math 4451, Spring 2026}
\fancyfoot[RE,RO]{\thepage}
\IfFileExists{upquote.sty}{\usepackage{upquote}}{}
\begin{document}
\begin{flushright}
\begin{minipage}{.33\textwidth}
\rightline{YOUR NAME}
\rightline{\href{mailto:YOUR EMAIL}{YOUR EMAIL}}
\rightline{\today}
\end{minipage}
\end{flushright}

\begin{center}
{\large{\textbf{Homework 5}}}
\end{center}

\begin{enumerate}

    \item (4 points) Beta-Binomial Conjugacy.
    \begin{enumerate}
      \item (3 points) Show that the Beta family of distributions is conjugate to the Binomial family of distributions (i.e., if the likelihood is Binomial$(n, p)$, then if the prior is $P \sim \text{Beta}(\alpha, \beta)$, then the posterior $P | X$ is a beta-distribution).
      \item (1 point) Use the result above to argue why a Uniform$(0, 1)$ prior for $P$ results in a Beta-distriubted posterior.
    \end{enumerate}

    \item (6 points) Approximate the following integrals using some form of Monte-Carlo. Describe the distribution used and the overall approach you picked.
    \begin{enumerate}
      \item (2 points) $\int_0^1 \sin\big(e^{-x}\big)x^2 \, dx$
      \item (2 points) $\int_0^2 e^{x^2}\, dx$
      \item (2 points) $\int_0^\infty \frac{e^{-x}}{1 + x^2}\, dx$
    \end{enumerate}

    \item (6 points) In class, we calculated the posterior mean $\Lambda | X$, where $X \sim \text{Pois}(\lambda)$, and the prior distribution was $\Lambda \sim U(0, 100)$. This required calculating the following integral:
    $$
    E[\Lambda | X = x^*] = \frac{\int_0^{100} \tau \cdot \tau^{\sum_i x_i^*}e^{-n\tau}\, d\tau}{\int_0^{100} \tau^{\sum_i x_i^*}e^{-n\tau}\, d\tau}
    $$
    (recall the denominator is just a normalizing constant.)
    To perform this calculation, it required making some transformations to avoid numeric overflow / underflow issues. In particular, we showed
    $$
    E[\Lambda | X = x^*] = \frac{\int_0^{100} \tau \cdot \tau^{\sum_i x_i^*}e^{-n\tau}\, d\tau}{\int_0^{100} \tau^{\sum_i x_i^*}e^{-n\tau}\, d\tau} = \frac{\int_0^{100} \lambda \exp\big\{g(\lambda) - g(\bar{x}^*)\big\}\, d\lambda}{\int_0^{100} \exp\big\{g(\lambda) - g(\bar{x}^*)\big\}\, d\lambda},
    $$
    where $g(\lambda) = \sum_i x_i^* \log \lambda - n\lambda$.

    \begin{enumerate}
      \item (3 points) In class, we used the \texttt{integrate} function in R to estimate the posterior mean. \textbf{Now, estimate the posterior mean using a Monte-Carlo method.} (The numerically-stable version will probably be easier).
      \item (3 points) Approximate the posterior variance $\text{Var}(\Lambda | X = x^*)$ using Monte-Carlo (HINT: try finding $E[\Lambda^2 | X = x^*]$).

    \end{enumerate}


\end{enumerate}


\end{document}
